\section{The Alembic Framework}
\label{sec:system}

\begin{figure}[t]
  \centering
  \includegraphics[width=0.78\columnwidth]{the_alembic.png}
  \caption{Structure of alembic: five LLM stages are
  interleaved with four deterministic gates that distill a raw
  repository, top to bottom, into a served MCP tool. Each stage runs inside a
  reset-from-checkpoint loop, and a structured \texttt{plan.json} contract flows
  to every downstream stage.}
  \label{fig:pipeline}
\end{figure}

Alembic is built on Google's Agent Development Kit and LiteLLM, so every agent
shares one configurable model. Its organizing principle is that \emph{the LLM
proposes and code disposes}: agents do the creative work (understanding the
repository, choosing what to wrap, writing tool bodies, diagnosing failures)
while every mechanical, correctness-bearing decision is moved out of the prompt
into a deterministic gate enforced by code, and stages hand off through a small
typed \texttt{plan.json} contract rather than free-text markdown. This keeps the
pipeline model-agnostic: a weaker model cannot slip an unchecked artifact past a
gate, and a stronger one is not slowed by hand-holding it does not need. Each
stage also runs inside a reset loop: on a gate failure the files it owns roll
back to a checkpoint and the stage reruns with a short note on what went wrong.

\subsection{Pipeline}
The pipeline (Figure~\ref{fig:pipeline}) alternates five agents with four gates.
The \emph{Explorer} clones the repository and proposes tool candidates as a typed
block (target symbol, sample and held-out arguments, a returns schema, and the
evidence grounding each test); a plan gate verifies every target against an AST
symbol table, drops hallucinated symbols, injects real parameter names, and fixes
the venv layout, producing the authoritative \texttt{plan.json}. The
\emph{Environment} agent builds the venv(s) with \texttt{uv} and \texttt{apt-get}
and downloads weights, recording every successful command into a reproducible
\texttt{setup.sh}; an environment gate runs a compatibility check and a
repo-import smoke test, so the Coder never writes against a broken venv.

Rather than emit the server directly, the \emph{Coder} writes each tool as a
plain, self-contained function with two kinds of tests (quick smoke checks and
evidence-based invocation checks), and an artefact gate statically confirms
each function compiles, imports, and carries its tests before any live run. The
\emph{Validator} invokes every tool, runs the suite, and batches failures into
one debugger call with problem-summary memory to keep fixes from oscillating;
each tool earns a two-level verdict: \emph{passed} when it runs and its unit
tests pass, \emph{perfect} when it also passes live invocation. Finally the
\emph{Wrapper} renders the FastMCP server deterministically from the verified
signatures (the LLM runs only if it fails to compile), after which the container
is committed to a versioned image served over MCP, exporting the tool functions,
tests, and \texttt{setup.sh} for reuse.

\subsection{Always-split environments}
The served server is two-layer: a FastMCP process on a modern Python (3.11+,
required by the MCP libraries) dispatches each call to a helper subprocess
running in the repository's own environment. Alembic therefore always keeps the
repository environment separate from the server environment, letting an older,
unmaintained repository run against the exact interpreter and pinned dependencies
its code was written for while the server layer stays current, avoiding the
package conflicts a single forced interpreter would introduce.
